\documentclass[proposal]{unnes-ta_plus}

% ---- metadata ----
\unnesSetTitle{Judul Proposal Skripsi Anda}
\unnesSetDegree{Sarjana \dots} % isi sesuai prodi
\unnesSetAuthor{Nama Mahasiswa}
\unnesSetNIM{1234567890}
\unnesSetProdi{Teknik Informatika}
\unnesSetFakultas{Fakultas Matematika dan Ilmu Pengetahuan Alam}
\unnesSetCity{Semarang}
\unnesSetYear{2026}
\unnesSetPembimbing{Nama Pembimbing}{NIP Pembimbing}

\begin{document}
\unnesFrontMatter

% Bagian awal proposal (panduan: sampul, halaman judul, persetujuan, daftar isi, dst.)
\maketitleunnes
\unnesPersetujuanPembimbing
\tableofcontents
\clearpage

\unnesMainMatter

\chapter{Pendahuluan}
\section{Latar Belakang}
\section{Rumusan Masalah}
\section{Tujuan Penelitian}
\section{Manfaat Penelitian}

\chapter{Kajian Pustaka}
\section{Tinjauan Pustaka}
\section{Landasan Teoretik}

\chapter{Metode Penelitian}
\section{Pendekatan, Jenis, dan Prosedur}
\section{Lokasi dan Waktu}
\section{Data dan Sumber Data}
\section{Teknik Analisis Data}

\chapter*{Daftar Pustaka}
\addcontentsline{toc}{chapter}{Daftar Pustaka}
% pakai biblatex/biber atau bibtex sesuai gaya prodi

\end{document}
